\documentclass[12pt, a4paper]{article}

% ── Core packages ─────────────────────────────────────────────────────────────
\usepackage[a4paper, top=2.5cm, bottom=2.5cm, left=3cm, right=3cm]{geometry}
\usepackage{amsmath, amssymb, amsthm}
\usepackage{mathtools}
\usepackage{bm}
\usepackage{booktabs}
\usepackage{array}
\usepackage{multirow}
\usepackage{xcolor}
\usepackage{hyperref}
\usepackage{cleveref}
\usepackage{graphicx}
\usepackage{float}
\usepackage{caption}
\usepackage{subcaption}
\usepackage{enumitem}
\usepackage{parskip}
\usepackage[protrusion=true,expansion=false]{microtype}
\usepackage{natbib}
\usepackage{tcolorbox}
\tcbuselibrary{breakable, skins, listings}
\usepackage{listings}
\usepackage[ruled, vlined, linesnumbered]{algorithm2e}
\usepackage{setspace}
\usepackage{titlesec}
\usepackage{abstract}
\usepackage{lipsum}

% ── Colour palette ────────────────────────────────────────────────────────────
\definecolor{darkblue}{RGB}{0,51,102}
\definecolor{midblue}{RGB}{0,102,179}
\definecolor{lightblue}{RGB}{219,234,248}
\definecolor{warmgrey}{RGB}{248,248,245}
\definecolor{codegrey}{RGB}{240,240,240}
\definecolor{darkred}{RGB}{150,20,20}
\definecolor{darkgreen}{RGB}{20,100,20}
\definecolor{pyblue}{RGB}{40,80,160}
\definecolor{pygreen}{RGB}{0,128,0}
\definecolor{pygrey}{RGB}{128,128,128}
\definecolor{pyorange}{RGB}{200,100,0}

% ── Hyperref ──────────────────────────────────────────────────────────────────
\hypersetup{
    colorlinks=true,
    linkcolor=darkblue,
    citecolor=midblue,
    urlcolor=midblue,
    pdftitle={A Unified Thermomechanical Model for MMW Drilling},
}

% ── Custom tcolorboxes ────────────────────────────────────────────────────────
\newtcolorbox{systembox}[1][]{
    enhanced, breakable,
    colback=lightblue!35,
    colframe=midblue,
    arc=3pt, boxrule=0.8pt,
    left=8pt, right=8pt, top=5pt, bottom=5pt,
    title={#1},
    fonttitle=\bfseries\color{darkblue},
}

\newtcolorbox{notebox}{
    enhanced, breakable,
    colback=warmgrey,
    colframe=darkblue!50,
    arc=2pt, boxrule=0.6pt,
    left=7pt, right=7pt, top=3pt, bottom=3pt,
}

\newtcolorbox{remarbox}{
    enhanced, breakable,
    colback=orange!8,
    colframe=pyorange!70,
    arc=2pt, boxrule=0.6pt,
    left=7pt, right=7pt, top=3pt, bottom=3pt,
}

% ── Python listings style ─────────────────────────────────────────────────────
\lstdefinestyle{pythonstyle}{
    language=Python,
    backgroundcolor=\color{codegrey},
    basicstyle=\ttfamily\footnotesize,
    keywordstyle=\color{pyblue}\bfseries,
    commentstyle=\color{pygrey}\itshape,
    stringstyle=\color{pyorange},
    numberstyle=\tiny\color{pygrey},
    numbers=left,
    stepnumber=1,
    numbersep=6pt,
    breaklines=true,
    showstringspaces=false,
    frame=single,
    framesep=4pt,
    rulecolor=\color{darkblue!30},
    tabsize=4,
    morekeywords={True, False, None},
}

% ── Math macros ───────────────────────────────────────────────────────────────
\newcommand{\pd}[2]{\dfrac{\partial #1}{\partial #2}}
\newcommand{\pdd}[2]{\dfrac{\partial^2 #1}{\partial #2^2}}
\newcommand{\vv}[1]{\bm{#1}}
\newcommand{\tens}[1]{\bm{\mathsf{#1}}}
\newcommand{\jump}[1]{\llbracket #1 \rrbracket}
\usepackage{stmaryrd}   % for \llbracket

% ── Section formatting ────────────────────────────────────────────────────────
\titleformat{\section}
    {\large\bfseries\color{darkblue}}{\thesection}{1em}{}
    [\vspace{-4pt}\color{midblue}\rule{\linewidth}{0.6pt}]
\titleformat{\subsection}
    {\normalsize\bfseries\color{darkblue!90}}{\thesubsection}{1em}{}
\titleformat{\subsubsection}
    {\normalsize\itshape\color{darkblue!80}}{\thesubsubsection}{1em}{}

\titlespacing*{\section}{0pt}{14pt}{6pt}
\titlespacing*{\subsection}{0pt}{10pt}{4pt}

% ── Algorithm style ───────────────────────────────────────────────────────────
\SetAlFnt{\small}
\SetAlCapFnt{\small\bfseries}
\SetAlCapNameFnt{\small\bfseries}
\newcommand\mycommfont[1]{\footnotesize\ttfamily\textcolor{pygrey}{#1}}
\SetCommentSty{mycommfont}

% ── Line spacing ──────────────────────────────────────────────────────────────
\setstretch{1.15}

% ─────────────────────────────────────────────────────────────────────────────
\begin{document}

% ── Title page ────────────────────────────────────────────────────────────────
\begin{titlepage}
\begin{center}
    \vspace*{1.5cm}
    {\Huge\bfseries\color{darkblue}
        A Unified Thermomechanical Model\\[8pt]
        for Millimetre-Wave Drilling:\\[8pt]
        Spallation and Vaporisation Regimes}\\[1.5cm]

    {\large\color{darkblue!70}\textit{Research Proposal}}\\[0.6cm]
    {\normalsize\color{black!70} \today}\\[2cm]

    \noindent\color{midblue}\rule{0.85\linewidth}{0.8pt}\\[0.4cm]
    \begin{minipage}{0.85\linewidth}
    \footnotesize\color{black!70}
    \textbf{Abstract.}
    Millimetre-wave (MMW) drilling is a contact-less excavation technology with
    potential to access deep geothermal resources currently beyond the reach of
    conventional mechanical drilling.  Existing numerical models treat either
    thermal spallation or vaporisation in isolation; however, both mechanisms
    operate simultaneously and their relative dominance shifts with depth and
    applied beam power.  This proposal describes a unified thermomechanical
    simulation framework that captures both material removal modes within a
    single fast engineering model.  The framework couples a Gaussian MMW beam
    model with a heterogeneous Voronoi microstructure, a pressure-dependent
    equation of state, a Drucker--Prager damage criterion, and a two-regime
    surface recession model.  Four research directions are addressed: (i)
    development of the cohesive multi-mode removal model; (ii) performance
    comparison across granite, basalt, and limestone; (iii) validation of the
    proposed spall-initiation time scaling law; and (iv) investigation of the
    depth-dependent regime transition driven by the equation of state.  The
    proposed work fills a significant gap in the MMW drilling literature and
    provides engineering design tools for geothermal well construction.
    \end{minipage}\\[0.4cm]
    \noindent\color{midblue}\rule{0.85\linewidth}{0.8pt}
\end{center}
\end{titlepage}

\tableofcontents
\newpage

% =============================================================================
\section{Introduction}
\label{sec:intro}
% =============================================================================

Accessing deep geothermal energy — at depths exceeding 5--10\,km in crystalline
basement rock — requires fundamentally new drilling technologies.  Conventional
rotary drilling faces exponentially increasing costs beyond 5\,km due to bit wear
in hard rock, tool failure under extreme temperatures, and slow penetration rates
\citep{tester2006future, augustine2009hydrothermal}.  Millimetre-wave (MMW) drilling
avoids these limitations entirely by delivering focused electromagnetic energy to
the rock face through a corrugated waveguide, fragmenting rock without mechanical
contact \citep{oglesby2014mmw, houde2021mmw, zhang2023thermal}.

Two fundamentally different material removal mechanisms operate in MMW drilling.
At moderate beam powers and shallow depths, \textbf{thermal spallation} dominates:
rapid surface heating generates large compressive thermal stresses that drive
pre-existing micro-cracks to coalesce, ejecting thin flakes from the surface well
below the melting point \citep{rauenzahn1989rock, kant2017theory}.  At higher
powers or greater depths — where confining pressure suppresses spall ejection —
\textbf{vaporisation} dominates: the rock surface melts and evaporates directly,
forming a melt pool that is expelled by the beam pressure \citep{zhang2023thermal}.
The transition between these regimes, and its dependence on material properties,
depth, and beam parameters, is poorly understood.

Existing numerical models address each regime in isolation.  The thermal model of
\citet{zhang2023thermal} captures vaporisation with high fidelity but does not
include the thermoelastic damage mechanics required to predict spallation.
Conversely, the damage models of \citet{vogler2020simulation} and
\citet{liu2023failure} capture grain-scale spallation but are not coupled to a
physically-based EOS or vaporisation model.  No current model treats both regimes
within a single unified framework, nor investigates the regime transition as a
function of depth.

\subsection*{Objectives and contributions}

This work proposes a unified fast-engineering simulation framework with four
specific research directions:

\begin{enumerate}[leftmargin=*, noitemsep]
    \item \textbf{Unified two-mode removal model.}  Develop a cohesive
    thermomechanical framework that captures both spallation and vaporisation
    within a single set of governing equations, with a physically based switching
    criterion between modes.

    \item \textbf{Material performance comparison.}  Apply the model to granite,
    basalt, and limestone to quantify the role of mineralogy and microstructure on
    spallation susceptibility, vaporisation threshold, and rate of penetration.

    \item \textbf{Spall initiation time scaling.}  Derive and validate an
    analytical scaling law for spall onset time as a function of beam power, rock
    thermal diffusivity, and microstructure; compare against available experimental
    and simulation data.

    \item \textbf{Depth-dependent regime transition.}  Couple the model to a
    pressure-dependent equation of state for rock and investigate how the
    spallation--vaporisation transition shifts with drilling depth.
\end{enumerate}

\begin{remarbox}
\textbf{Physics and numerics added beyond the baseline framework.}
In addition to the standard thermal--mechanical--damage coupling, the present
model incorporates:
(a) \textbf{Enthalpy-based conservation} replacing the temperature PDE, handling
phase transitions without a modified specific heat;
(b) \textbf{Beer--Lambert attenuation} with frequency-dependent $\alpha \propto f$
linking gyrotron selection to removal regime;
(c) \textbf{Mie--Gr\"{u}neisen equation of state} for pressure-dependent
elastic modulus at depth;
(d) \textbf{Drucker--Prager} failure criterion accounting for pressure-dependent
rock strength;
(e) \textbf{Voronoi microstructure} with grain-boundary cohesive zones as the
physical mechanism driving spallation;
(f) \textbf{temperature-dependent fracture toughness} $K_{Ic}(T)$;
(g) a \textbf{confining pressure correction} to the damage threshold; and
(h) \textbf{block-structured AMR} via AMReX/ALAMO with two refinement levels.
\end{remarbox}

% =============================================================================
\section{Physical Framework}
\label{sec:physics}
% =============================================================================

\subsection{MMW energy deposition and the Beer--Lambert law}
\label{sec:beam}

\subsubsection{Beam profile and volumetric absorption}

The MMW power deposited per unit volume follows the Beer--Lambert law.  For a
beam incident on the rock surface at $z = z_0$, the volumetric source is
\begin{equation}
    \dot{Q}(\vv{x}, T) = P(\vv{x})\,\alpha(T)\,
                          \exp\!\left(-\alpha(T)(z_0 - z)\right),
    \label{eq:qdot}
\end{equation}
where $\alpha(T)$ is the temperature-dependent attenuation coefficient
(units m$^{-1}$) and $P(\vv{x})$ is the surface power intensity.
For a Gaussian beam centred at $(x_0, y_0)$,
\begin{equation}
    P_\mathrm{GB}(x,y,z) = \frac{2P_0}{\pi\omega_0^2}
            \left(\frac{\omega_0}{\omega(z)}\right)^{\!2}
            \exp\!\left(\frac{-2r^2}{\omega^2(z)}\right),
    \qquad
    r = \sqrt{(x-x_0)^2 + (y-y_0)^2},
    \label{eq:intensity}
\end{equation}
with spot size
\begin{equation}
    \omega(z) = \omega_0\sqrt{1 + \left(\frac{\lambda\,(z_0 - z)}
                                         {\pi\omega_0^2}\right)^{\!2}},
    \label{eq:waist}
\end{equation}
where $P_0$ is beam power, $\omega_0$ the beam waist at the waveguide exit,
and $\lambda$ the MMW wavelength.

\subsubsection{Role of the attenuation coefficient}
\label{sec:attenuation}

The attenuation coefficient $\alpha$ is the single most important parameter
governing the \emph{spatial shape} of energy deposition, and therefore
exerts a first-order influence on whether spallation or vaporisation is
the dominant removal mode.

\paragraph{Physical definition.}
For a dielectric rock material, the attenuation coefficient is related to
the imaginary part of the complex permittivity
$\varepsilon^* = \varepsilon' - i\varepsilon''$ by
\begin{equation}
    \alpha = \frac{2\omega}{c}\sqrt{
        \frac{\mu_r}{2}\!\left(\sqrt{{\varepsilon'}^2 + {\varepsilon''}^2}
                               - \varepsilon'\right)},
    \label{eq:alpha_def}
\end{equation}
where $\omega = 2\pi c/\lambda$ is the angular frequency and
$\mu_r \approx 1$.  In the low-loss limit ($\varepsilon'' \ll \varepsilon'$),
this simplifies to
\begin{equation}
    \alpha \approx \frac{\pi}{\lambda}\,
                   \frac{\varepsilon''}{\sqrt{\varepsilon'}}
           = \frac{\pi\,\tan\delta}{\lambda}\sqrt{\varepsilon'},
    \label{eq:alpha_approx}
\end{equation}
where $\tan\delta = \varepsilon''/\varepsilon'$ is the loss tangent.
Two key consequences follow:
\begin{enumerate}[noitemsep]
    \item \textbf{$\alpha \propto 1/\lambda \propto f$}: the attenuation
    coefficient is proportional to the gyrotron operating frequency.
    A 28\,GHz gyrotron ($\lambda \approx 10.7\,\mathrm{mm}$) produces
    roughly $3\times$ less surface attenuation than a 94\,GHz system
    ($\lambda \approx 3.2\,\mathrm{mm}$) for the same rock
    \citep{woskov2009mmw}.

    \item \textbf{Temperature feedback:} the loss tangent $\tan\delta(T)$
    increases as minerals approach their melting point, so $\alpha$ rises
    with temperature, further concentrating energy deposition at the
    surface — a positive feedback that accelerates phase transitions.
\end{enumerate}

\paragraph{Effect on thermal gradient and removal regime.}
The characteristic penetration depth is $\ell_\alpha = 1/\alpha$.
The induced surface thermal gradient scales as
\begin{equation}
    \left.\frac{\partial T}{\partial z}\right|_{z=z_0}
    \;\sim\; \frac{P\,\alpha}{\kappa}
    \;=\; \frac{P}{\kappa\,\ell_\alpha},
    \label{eq:grad_scale}
\end{equation}
giving rise to two limiting regimes:
\begin{itemize}[noitemsep]
    \item \textbf{High $\alpha$ (high frequency, $\ell_\alpha \ll
    h_\mathrm{spall}$):} energy deposits in a thin near-surface layer;
    steep gradients drive large compressive thermal stresses favouring
    \emph{spallation}.

    \item \textbf{Low $\alpha$ (low frequency, $\ell_\alpha \gg
    h_\mathrm{spall}$):} energy spreads through a thick rock volume;
    gradients are mild, stress concentrations at grain boundaries are
    reduced, and the surface must reach $T_v$ before removal, favouring
    \emph{vaporisation}.
\end{itemize}
The mode transition occurs when $\ell_\alpha \sim h_\mathrm{spall} \sim
C_h\sqrt{\alpha_\mathrm{th}\,t_\mathrm{spall}}$, directly linking the
gyrotron frequency to the dominant removal mechanism.

\paragraph{Proposed test: attenuation sensitivity study.}
\label{para:atten_test}
To quantify this, a controlled numerical experiment is proposed.  All
parameters are held fixed except the attenuation coefficient, which is
swept over
$\alpha \in \{10,\, 30,\, 100,\, 300,\, 1000\}\,\mathrm{m}^{-1}$,
corresponding approximately to gyrotron frequencies of 28, 45, 94, 140,
and 300\,GHz respectively for granitic rock \citep{woskov2009mmw}.
Each simulation records:
\begin{enumerate}[noitemsep]
    \item The depth profile $\partial T/\partial z(z)$ at $t = 1\,\mathrm{s}$
          to visualise the change in energy localisation;
    \item The onset-spallation time $t_\mathrm{spall}$, or ``vaporisation
          dominant'' if $T_v$ is reached first;
    \item The specific energy (J\,mm$^{-3}$ removed) for each case.
\end{enumerate}
The expected result is a monotone decrease in $t_\mathrm{spall}$ with
increasing $\alpha$, saturating when $\ell_\alpha$ drops below the mean
grain size (at which point grain-boundary heterogeneity, not the beam
profile, controls damage localisation).  A crossover from spallation to
vaporisation is expected at low $\alpha$, providing a direct experimental
prediction linking gyrotron selection to drilling performance.

\noindent The total source entering the enthalpy equation is
\begin{equation}
    q(\vv{x},T) = \dot{Q}(\vv{x},T)
                 - \varepsilon\sigma\bigl(T^4 - T_0^4\bigr)
                 - h\bigl(T - T_0\bigr),
    \label{eq:source}
\end{equation}
with the loss terms applied only at the ablation surface.

\subsection{Enthalpy conservation and phase tracking}
\label{sec:heat}

Rather than evolving temperature directly, the model tracks the volumetric
enthalpy $H$.  This formulation, introduced by \citet{alexiades2010enthalpy},
handles phase transitions naturally without the need for explicit front-tracking
or a modified specific heat.  Under the constant-pressure assumption ($\mathrm{d}p
= 0$), the volumetric enthalpy satisfies
\begin{equation}
    \mathrm{d}H = \rho\,c_p\,\mathrm{d}T,
    \label{eq:dH}
\end{equation}
and the conservation law takes the form
\begin{equation}
    \pd{H}{t}
    - \nabla\cdot\!\left[\frac{\kappa(T,D,\vv{x})}{\rho(T)\,c_p(T)}\,\nabla H\right]
    = q(\vv{x}, T),
    \label{eq:enthalpy_pde}
\end{equation}
where the source term $q$ is the net volumetric power input:
\begin{equation}
    q(\vv{x}, T) = \dot{Q}(\vv{x},T)
                   - \varepsilon\sigma\bigl(T^4 - T_0^4\bigr)
                   - h\bigl(T - T_0\bigr).
    \label{eq:source_q}
\end{equation}
The radiation and convection losses (the final two terms) are resolved only at
the ablation surface and are computed using the surface temperature $T_0$ and
the vertical extent of the corresponding surface cell $\mathrm{d}z$.

\subsubsection{Temperature--enthalpy inversion}

Since all material properties ($\kappa$, $\rho$, $c_p$, $E$, $\beta$) are
functions of temperature, the enthalpy field must be inverted to recover $T$
at every time step.  This inversion is governed by four phase-indicator
enthalpies \citep{alexiades2010enthalpy}:
\begin{itemize}[noitemsep]
    \item $H^S$: enthalpy at the onset of melting (solidus);
    \item $H^L = H^S + \Delta H^\mathrm{fus}$: enthalpy at the end of melting
          (liquidus);
    \item $H^{LV}$: enthalpy at the onset of vaporisation;
    \item $H^V = H^{LV} + \Delta H^\mathrm{vap}$: enthalpy at the completion
          of vaporisation.
\end{itemize}

\noindent\textbf{Solid phase} ($H < H^S$, $T < T_m$):
\begin{equation}
    H^S - H(T) = \int_{T}^{T_m} \rho(T')\,c_p(T')\,\mathrm{d}T'.
    \label{eq:H_solid}
\end{equation}

\noindent\textbf{Solid--liquid transition} ($H^S \le H \le H^L$):
temperature is constant at $T = T_m$ while phase fractions evolve.
Mixture properties are computed via linear combination of the
solid ($\Lambda^S$) and liquid ($\Lambda^L$) volume fractions:
\begin{equation}
    \Lambda^L = \frac{H - H^S}{\Delta H^\mathrm{fus}},
    \qquad
    \Lambda^S = 1 - \Lambda^L,
    \label{eq:phase_fracs}
\end{equation}
so that, for example, $\rho = \Lambda^L\rho^L + \Lambda^S\rho^S$.  The
same linear-combination rule is applied to $\kappa$, $c_p$, and $\beta$.

\noindent\textbf{Liquid phase} ($H^L < H < H^{LV}$): analogous to the solid
phase with fully liquid properties.

\noindent\textbf{Liquid--vapour transition} ($H^{LV} \le H \le H^V$):
temperature is constant at $T = T_v$; liquid and vapour volume fractions
are computed analogously to \cref{eq:phase_fracs} using $\Delta H^\mathrm{vap}$.

\noindent\textbf{Vapour phase} ($H > H^V$): material has been removed and the
surface $\phi = 0$ is advanced accordingly.

\noindent The spatial thermal conductivity field appearing in
\cref{eq:enthalpy_pde} is the heterogeneous, damage-modified field
$\kappa(T, D, \vv{x})$ defined in \S\ref{sec:voronoi}.

\subsection{Equation of state for elastic modulus}
\label{sec:eos}

At depth, rock is subject to significant confining pressure $P_c$.  The elastic
modulus is no longer a constant material property but depends on both pressure
and temperature (or internal energy $e$).  We adopt a \textbf{Mie--Gr\"{u}neisen}
form for the pressure--energy relation:
\begin{equation}
    P(V,e) = P_H(V) + \frac{\Gamma(V)}{V}\left(e - e_H(V)\right),
    \label{eq:mie-gruneisen}
\end{equation}
where $V = 1/\rho$ is the specific volume, $P_H$ and $e_H$ are the
Hugoniot reference pressure and internal energy, and $\Gamma(V)$ is the
Gr\"{u}neisen parameter.  For a linear $u_s$--$u_p$ Hugoniot,
\begin{equation}
    P_H(\eta) = \frac{\rho_0\,c_0^2\,\eta}{(1 - s\eta)^2}, \qquad
    \eta = 1 - \frac{\rho_0}{\rho},
    \label{eq:hugoniot}
\end{equation}
where $c_0$ is the ambient bulk sound speed and $s$ the Hugoniot slope
coefficient.  The pressure-dependent Young's modulus is then recovered as
\begin{equation}
    E(P, T) = \rho\,c_s^2(P)\,\frac{3(1-2\mu)}{1},
    \qquad
    c_s^2(P) = c_0^2 + \frac{\partial P}{\partial \rho}\bigg|_e,
    \label{eq:E_EOS}
\end{equation}
where $c_s$ is the pressure-modified longitudinal sound speed.  At low pressures
and temperatures this reduces to the standard elastic modulus $E_0$.  The
practical consequence for MMW drilling is that $E$ increases with depth,
raising the LRST and spallation threshold, and eventually rendering spallation
impossible — at which point vaporisation becomes the sole removal mechanism.

\subsection{Stress criterion: von Mises, Drucker--Prager, and Hoek--Brown}
\label{sec:stress_criterion}

The choice of failure criterion significantly affects which regions of the rock
are predicted to damage and which removal mode dominates.  We discuss three
candidates and justify our choice.

\subsubsection{Von Mises criterion}

The von Mises criterion,
\begin{equation}
    \sigma_v = \sqrt{\tfrac{1}{2}\!\left[
        (\sigma_1-\sigma_2)^2 + (\sigma_2-\sigma_3)^2 + (\sigma_3-\sigma_1)^2
        \right]} \ge \sigma_Y,
    \label{eq:vonmises}
\end{equation}
is \emph{pressure-independent}: it predicts the same yield stress regardless of
the mean stress $p = -(\sigma_1+\sigma_2+\sigma_3)/3$.  This is appropriate for
metals but fundamentally incorrect for rock, which is significantly stronger in
compression than in tension, and whose strength increases with confining pressure.
Under a MMW beam, the near-surface layer is in low confinement while material at
depth is under $P_c$.  Von Mises cannot capture this contrast and will overpredict
damage at depth.  It is retained in some thermal spallation studies
\citep{hu2018lrst} purely for simplicity.

\subsubsection{Drucker--Prager criterion (recommended)}

The Drucker--Prager (DP) criterion is the smooth, circular-cone generalisation
of Mohr--Coulomb in principal-stress space:
\begin{equation}
    F_\mathrm{DP} = \sqrt{J_2} + \alpha_\mathrm{DP}\,I_1 - k_\mathrm{DP} \ge 0,
    \label{eq:DP}
\end{equation}
where $I_1 = \sigma_{kk}$ is the first stress invariant (trace),
$J_2 = \tfrac{1}{2}s_{ij}s_{ij}$ is the second deviatoric stress invariant, and
the material constants are
\begin{equation}
    \alpha_\mathrm{DP} = \frac{2\sin\phi}{\sqrt{3}(3 - \sin\phi)},
    \qquad
    k_\mathrm{DP} = \frac{6c\cos\phi}{\sqrt{3}(3 - \sin\phi)},
    \label{eq:DP_params}
\end{equation}
with $\phi$ the internal friction angle and $c$ the cohesion.  DP correctly
predicts that rock is weaker in tension than compression, strengthens with
confinement, and — critically for deep drilling — the damage threshold increases
with $P_c = I_1/3$.  We adopt this criterion as the primary failure measure for
both tensile-dominated spallation and compressive shear failure at depth.

\subsubsection{Hoek--Brown spalling limit}

For the specific case of near-surface spallation under low confinement,
\citet{vogler2020simulation} showed that a bilinear spalling limit cut-off on
the Hoek--Brown envelope better captures the reduced strength of the
grain-boundary zone:
\begin{equation}
    \frac{\sigma_1}{\sigma_3} \ge \xi_s \approx 8,
    \label{eq:HB_spall}
\end{equation}
where $\xi_s$ is the spalling ratio.  This criterion is applied as a secondary
check for spall \emph{detachment} after the DP criterion indicates that $D \to 1$
in a connected near-surface zone.

\subsubsection{Practical implementation}

In the unified model, failure is assessed hierarchically:
\begin{enumerate}[noitemsep]
    \item Evaluate $F_\mathrm{DP}$; if $F_\mathrm{DP} \ge 0$ update damage $D$
          via \cref{eq:damage_evo}.
    \item If $D \to 1$ over a connected zone \emph{and} $\sigma_1/\sigma_3
          \ge \xi_s$, trigger spall detachment.
    \item If surface temperature $T \ge T_v$ (vaporisation threshold), activate
          the vaporisation removal mode regardless of mechanical criteria.
\end{enumerate}

\subsection{Voronoi microstructure and grain-boundary physics}
\label{sec:voronoi}

Grain-scale heterogeneity is not a refinement — it is the physical mechanism
responsible for spallation.  \citet{vogler2020simulation} and
\citet{liu2023failure} both demonstrated that \emph{homogeneous} material models
predict zero spallation damage under conditions where experiments show abundant
fracture.  The Voronoi tessellation is therefore central to the framework.

\subsubsection{Tessellation and phase assignment}

A Voronoi tessellation of the domain $\Omega$ is constructed from $N$ seed points
$\{\vv{x}_k\}_{k=1}^N$, chosen so that the average grain radius matches target
mineralogy:
\begin{equation}
    R_a = \sum_{i=1}^m \phi_i\,r_i, \qquad
    H = \sqrt{\sum_{i=1}^m \phi_i\!\left(\frac{r_i}{R_a}-1\right)^{\!2}},
    \label{eq:grain_size}
\end{equation}
where $\phi_i$ and $r_i$ are the volume fraction and average radius of mineral
phase $i$, and $H$ is the dimensionless heterogeneity index \citep{liu2023failure}.
Each Voronoi cell is assigned a mineral phase by drawing from the multinomial
distribution $\{\phi_k\}$.  The resulting piecewise-constant property fields
$\kappa(\vv{x})$, $\beta(\vv{x})$, $E(\vv{x})$ carry discontinuities at every
grain boundary.

\subsubsection{Why heterogeneity drives spallation}

Three grain-scale mechanisms generate the stress concentrations required for
damage nucleation:
\begin{enumerate}[noitemsep]
    \item \textbf{Differential thermal conduction.}  Quartz
    ($\kappa \approx 7.7\,\mathrm{W\,m^{-1}K^{-1}}$) conducts heat $3\times$
    faster than feldspar ($\kappa \approx 2.3\,\mathrm{W\,m^{-1}K^{-1}}$), so
    adjacent grains reach different temperatures under the same heat flux,
    generating a $\nabla T$ discontinuity at their shared boundary.

    \item \textbf{Thermal expansion mismatch.}  Quartz
    ($\beta \approx 1.67\times10^{-5}\,\mathrm{K}^{-1}$) expands $\approx
    4\times$ more than K-feldspar ($\beta \approx 0.37\times10^{-5}\,\mathrm{K}^{-1}$).
    The resulting intergranular stress concentration at a quartz--feldspar
    boundary scales as
    \begin{equation}
        \Delta\sigma \sim
        \frac{E_1\,E_2}{E_1+E_2}\,(\beta_1 - \beta_2)\,\Delta T.
        \label{eq:grain_stress}
    \end{equation}
    For $\Delta T \approx 200\,\mathrm{K}$, this yields $\Delta\sigma
    \approx 20$--$50\,\mathrm{MPa}$, sufficient to exceed grain-boundary
    tensile strength.

    \item \textbf{Quartz $\alpha$--$\beta$ phase transition.}  At
    $573\,^\circ\mathrm{C}$ (at atmospheric pressure), quartz undergoes a
    displacive phase transition from trigonal to hexagonal symmetry with a
    volumetric expansion of $\approx 0.63\%$, with anisotropic linear strains
    of $0.24\%$ and $0.14\%$ along the principal $a$ and $c$ axes respectively
    \citep{carpenter1998quartz, johnson2021quartz}.  The transition temperature
    is pressure-dependent, increasing at $\approx
    0.255\,^\circ\mathrm{C\,MPa}^{-1}$ \citep{shen1993quartz}, so at
    1\,km drilling depth ($\sim 25$\,MPa) the transition shifts to
    $\approx 580\,^\circ\mathrm{C}$ --- still within the thermal spallation
    regime.  At the transition, the thermal expansion coefficient diverges
    before dropping sharply, and Young's modulus falls to a local minimum
    \citep{haussuehl1960quartz}.  This produces a sharp spike in the effective
    $\beta_\mathrm{qtz}(T)$, a pulse of intergranular stress independent of
    the background thermal gradient, and a strong burst of acoustic-emission
    activity confirming micro-crack nucleation at grain boundaries
    \citep{griffin1995acoustic}.
\end{enumerate}

\subsubsection{Grain-boundary cohesive zones}

Grain boundaries are modelled as zero-thickness cohesive interface elements with
a bilinear traction-separation law:
\begin{equation}
    t_n(\delta_n) =
    \begin{cases}
        K_n\,\delta_n & 0 \le \delta_n < \delta_c \\
        K_n\,\delta_c\,\dfrac{\delta_\mathrm{max} - \delta_n}
                              {\delta_\mathrm{max} - \delta_c}
        & \delta_c \le \delta_n < \delta_\mathrm{max} \\
        0 & \delta_n \ge \delta_\mathrm{max}
    \end{cases}
    \label{eq:czm}
\end{equation}
where $K_n$ is the interface stiffness, $\delta_c$ the crack-opening displacement
at peak traction, and $\delta_\mathrm{max}$ the failure opening.  The grain
boundary tensile strength is set to $f_{t0}^\mathrm{gb} \approx 0.3$--$0.5\,
f_{t0}^\mathrm{grain}$ and the interface fracture energy
$G_c^\mathrm{gb} = K_n\,\delta_\mathrm{max}^2/2$ is calibrated to the mode-I
fracture toughness of grain boundaries in granite.  As damage accumulates ($D
\to 1$ at the interface), the thermal conductance also degrades:
\begin{equation}
    h_\mathrm{gb}(D) = h_{\mathrm{gb},0}\,(1-D),
    \label{eq:gb_conductance}
\end{equation}
creating an insulating crack that amplifies local temperature gradients and
accelerates further damage — the positive feedback loop responsible for spall
detachment.

\subsubsection{Homogenised vs heterogeneous conductivity}

Within each grain, the damage-modified conductivity is
\begin{equation}
    \kappa(\vv{x},T,D) = \kappa_k(T)\,\exp(-\alpha_D\,D),
    \label{eq:kappa_damage}
\end{equation}
and the bulk effective value is computed via arithmetic-harmonic averaging over
mineral fractions \citep{vogler2020simulation}:
\begin{equation}
    \kappa_\mathrm{eff} = \frac{1}{2}\!\left(
        \sum_k \phi_k\,\kappa_k
        + \frac{1}{\displaystyle\sum_k \phi_k/\kappa_k}
    \right).
    \label{eq:keff}
\end{equation}

\subsection{Damage evolution}
\label{sec:damage}

The damage variable $D \in [0,1]$ evolves continuously under the Drucker--Prager
overstress.  The damage driving force is the normalised excess stress
\begin{equation}
    \Psi = \frac{F_\mathrm{DP}}{k_\mathrm{DP}},
    \label{eq:psi}
\end{equation}
and the evolution law is
\begin{equation}
    \pd{D}{t} =
    A_D(\vv{x})\left(\frac{\Psi}{1 + \Psi_0}\right)^{\!n}(1-D)^m
    \cdot\mathbf{1}[\Psi \ge \Psi_0] \cdot \mathbf{1}[\dot{D} \ge 0],
    \label{eq:damage_evo}
\end{equation}
where $\Psi_0$ is the initiation threshold (approximately $\Psi_0 \approx 0.15$
for granite, calibrated against $f_{b,0} \approx 0.82$ from \citealt{hu2018lrst}),
$n$ controls the acceleration rate, and $(1-D)^m$ ensures saturation.  The rate
coefficient $A_D(\vv{x})$ is mineral-phase specific and elevated at grain
boundaries.  Stiffness degrades as
\begin{equation}
    E(\vv{x},D) = E(\vv{x},P,T)\,(1-D), \qquad G(\vv{x},D) = G_0(\vv{x})\,(1-D).
    \label{eq:stiffness_degrade}
\end{equation}

\subsubsection{Confining pressure correction at depth}

At drilling depth $z_d$, the lithostatic confining pressure is
$P_c = \rho_r\,g\,z_d$, where $\rho_r$ is the average overburden density.
This shifts the DP yield surface outward, increasing the effective initiation
threshold:
\begin{equation}
    \Psi_0(P_c) = \Psi_0^{(0)}\left(1 + \frac{\alpha_\mathrm{DP}\,P_c}{k_\mathrm{DP}}\right).
    \label{eq:depth_correction}
\end{equation}
Physically, greater confining pressure suppresses spall ejection and raises the
energy required to initiate fracture, driving the system toward the vaporisation
regime.

\subsubsection{Temperature-dependent fracture toughness}

Experimental data \citep{rossi2018effects} show that $K_{Ic}$ decreases
approximately linearly with temperature for granitic rocks:
\begin{equation}
    K_{Ic}(T) = K_{Ic}^{(0)}\!\left(1 - \frac{T - T_\mathrm{ref}}
                                               {T_\mathrm{melt}}\right), \quad
    T \in [T_\mathrm{ref},\, T_\mathrm{melt}].
    \label{eq:Kic_T}
\end{equation}
This is incorporated via the grain-boundary cohesive energy
$G_c(T) \propto K_{Ic}^2(T)/E(T)$, which decreases with temperature and
promotes crack propagation at lower stress levels.

\subsection{Spall thickness scaling and rate of penetration}
\label{sec:spall_scaling}

Once a connected damage zone forms and the spall detachment criterion is
satisfied, the spall thickness is estimated from the thermal diffusion length at
the onset time:
\begin{equation}
    h_\mathrm{spall} = C_h\sqrt{\alpha\,t_\mathrm{spall}},
    \qquad
    \alpha = \frac{\kappa_\mathrm{eff}}{\rho\,C_p},
    \label{eq:spall_thickness}
\end{equation}
where $C_h$ is a geometry factor of order unity.  The onset time is estimated
analytically from the heat-flux solution \citep{hu2018lrst, carslaw1959},
\begin{equation}
    Q = \frac{A\,K_r}{2(\alpha/\pi)^{1/2}},
    \label{eq:Q_onset}
\end{equation}
where $A = \mathrm{d}T_s/\mathrm{d}t^{1/2}$ is the slope of surface temperature
against $\sqrt{t}$.  Inverting for $t_\mathrm{spall}(Q)$ and substituting into
\cref{eq:spall_thickness} gives the scaling
\begin{equation}
    h_\mathrm{spall} \propto C_h\,\frac{K_r}{\rho\,C_p}\,
    \left(\frac{2K_r}{Q}\right)^2 \frac{1}{\sqrt{\pi\alpha}}.
    \label{eq:spall_scaling_explicit}
\end{equation}
The \textbf{rate of penetration} from spallation is
\begin{equation}
    \mathrm{RoP}_\mathrm{spall} = \frac{h_\mathrm{spall}}{t_\mathrm{spall}},
    \label{eq:rop_spall}
\end{equation}
and from vaporisation it is
\begin{equation}
    \mathrm{RoP}_\mathrm{vap} = \frac{(1-R)\,P_0}
                                      {\rho\,L_\mathrm{tot}\,A_\mathrm{beam}},
    \qquad
    L_\mathrm{tot} = C_p(T_v - T_0) + L_m + L_v,
    \label{eq:rop_vap}
\end{equation}
where $A_\mathrm{beam}$ is the beam footprint area.  The overall RoP is
\begin{equation}
    \mathrm{RoP} = \max(\mathrm{RoP}_\mathrm{spall},\;\mathrm{RoP}_\mathrm{vap}),
    \label{eq:rop_total}
\end{equation}
with the dominant mode determined by whichever criterion is first satisfied at
each time step.

% =============================================================================
\section{Numerical Model}
\label{sec:numerics}
% =============================================================================

\subsection{Overall architecture}

The simulation is structured as a staggered operator-split loop over time steps.
At each step, four sequential solves are performed:
(1) thermal solve for $T^{n+1}$; (2) mechanical solve for $\vv{u}^{n+1}$;
(3) damage update for $D^{n+1}$; (4) surface and regime update.
This decoupling allows each sub-problem to use its natural solver (implicit
diffusion, static equilibrium, explicit integration) without forming a
monolithic Jacobian.

\begin{algorithm}[H]
\caption{Main MMW drilling simulation loop (AMReX/ALAMO)}
\label{alg:main}
\DontPrintSemicolon
\SetKwInOut{Input}{Input}\SetKwInOut{Output}{Output}
\Input{Beam parameters $(P_0, \omega_0, \lambda, \alpha(T))$; rock properties;
       depth $z_d$; Voronoi microstructure; time span $[0, t_\mathrm{end}]$}
\Output{$H(\vv{x},t)$, $T(\vv{x},t)$, $D(\vv{x},t)$, surface profile $\phi(t)$,
        $\mathrm{RoP}(t)$}
\BlankLine
Initialise AMR hierarchy; base mesh $\Delta x^0$; 2 refinement levels\;
Initialise: $H^0 = H(T_\mathrm{amb})$, $D^0 = 0$, $\phi^0 = 0$, $z_0 = z_d$\;
Compute lithostatic pressure $P_c = \rho_r g z_d$\;
Compute EOS-adjusted moduli $E(\vv{x}, P_c, T^0)$ per grain\;
\BlankLine
\For{$n = 0, 1, 2, \ldots$ until $t^n \ge t_\mathrm{end}$}{
    \tcc{Regrid AMR hierarchy every $N_\mathrm{regrid}$ steps}
    \If{$n \bmod N_\mathrm{regrid} = 0$}{
        Tag cells by $|\nabla H|$, $|\nabla D|$, beam footprint\;
        Regenerate fine levels; propagate phase array by injection\;
    }
    \BlankLine
    \tcc{Step 1: Enthalpy solve (implicit Crank--Nicolson per AMR level)}
    Evaluate $\dot{Q}(\vv{x}, z_0 - z, \alpha(T^n))$ via \cref{eq:qdot}\;
    Assemble source $q^n$ via \cref{eq:source}\;
    Solve \cref{eq:enthalpy_pde} implicitly for $H^{n+1}$ at each level\;
    Perform reflux and average-down across levels\;
    \BlankLine
    \tcc{Step 1b: Enthalpy inversion to recover $T^{n+1}$}
    \For{each cell $\vv{x}_i$}{
        Invert $H^{n+1}(\vv{x}_i) \to T^{n+1}(\vv{x}_i)$ via
        phase indicators $\{H^S, H^L, H^{LV}, H^V\}$ (\S\ref{sec:heat})\;
        Update phase fractions $\Lambda^S$, $\Lambda^L$, $\Lambda^V$\;
        Update $\kappa$, $\rho$, $c_p$, $\beta$ from phase and Voronoi phase\;
    }
    \BlankLine
    \tcc{Step 2: Mechanical solve}
    Compute thermal strains $\varepsilon^{th} = \beta(\vv{x})\,(T^{n+1}-T_0)$\;
    Assemble degraded stiffness $\tens{K}(D^n, E(\vv{x}, P_c, T^{n+1}))$\;
    Solve for $\vv{u}^{n+1}$ and stresses $\sigma^{n+1}$\;
    \BlankLine
    \tcc{Step 3: Damage update}
    Evaluate DP criterion $F_\mathrm{DP}(\sigma^{n+1}, P_c)$\;
    \If{$F_\mathrm{DP} \ge 0$ at grain boundary cells}{
        Apply CZM traction-separation update \cref{eq:czm}\;
    }
    Integrate damage ODE \cref{eq:damage_evo} with RK4 for $D^{n+1}$\;
    Apply irreversibility: $D^{n+1} \leftarrow \max(D^{n+1},\, D^n)$\;
    Update $E \leftarrow E(1-D^{n+1})$\;
    \BlankLine
    \tcc{Step 4: Surface and regime update}
    \uIf{$D^{n+1} \to 1$ over connected zone \textbf{and}
          $\sigma_1/\sigma_3 \ge \xi_s$ \textbf{(spallation)}}{
        $h_s \leftarrow C_h\sqrt{\alpha_\mathrm{th}\,\Delta t_\mathrm{spall}}$\;
        $\phi \leftarrow \phi - h_s$, $z_0 \leftarrow z_0 - h_s$;
        reset $D=0$, $H = H(T_\mathrm{amb})$ in detached zone\;
        Record $\mathrm{RoP} \leftarrow h_s / \Delta t_\mathrm{spall}$\;
    }
    \uElseIf{$H^{n+1} \ge H^V$ at surface \textbf{(vaporisation)}}{
        $h_v \leftarrow \mathrm{RoP}_\mathrm{vap}\cdot\Delta t$
        via \cref{eq:rop_vap}\;
        $\phi \leftarrow \phi - h_v$, $z_0 \leftarrow z_0 - h_v$\;
        Record $\mathrm{RoP} \leftarrow h_v / \Delta t$\;
    }
    Advance time; adapt $\Delta t^\ell$ via \cref{eq:cfl_amr}\;
}
\end{algorithm}

\subsection{Voronoi mesh generation}
\label{sec:voronoi_numerics}

\begin{algorithm}[H]
\caption{Voronoi microstructure generation}
\label{alg:voronoi}
\DontPrintSemicolon
\SetKwInOut{Input}{Input}\SetKwInOut{Output}{Output}
\Input{Domain $\Omega = [L_x \times L_y \times L_z]$; target mineral fractions
       $\{\phi_k\}$; target grain radii $\{r_k\}$; FEM mesh nodes}
\Output{Grain-phase array \texttt{phase[node]};
        grain-boundary element flag \texttt{is\_gb[elem]}}
\BlankLine
Compute target seed count $N = |\Omega|/(\tfrac{4}{3}\pi R_a^3)$\;
Generate $N$ initial seeds via Poisson disk sampling in $\Omega$\;
\BlankLine
\tcc{Lloyd's algorithm: regularise grain sizes}
\For{$\ell = 1$ to $N_\mathrm{Lloyd}$ iterations}{
    Compute Voronoi cells $\{V_k\}$ from current seeds\;
    Move each seed to centroid of its cell\;
    \If{$|\mathrm{seed} - \mathrm{centroid}| < \epsilon_\mathrm{tol}$}{break}
}
\BlankLine
\tcc{Phase assignment}
Draw phase labels $p_k \sim \mathrm{Multinomial}(\{\phi_k\})$ for $k=1\ldots N$\;
Verify volume fractions match targets within 1\%\;
\BlankLine
\tcc{Map to FEM mesh}
\For{each FEM node $\vv{x}_i$}{
    $k^* \leftarrow \arg\min_k |\vv{x}_i - \mathrm{seed}_k|$\;
    \texttt{phase[$i$]} $\leftarrow p_{k^*}$\;
}
\BlankLine
\tcc{Identify grain boundary elements}
\For{each FEM element $e$}{
    \If{nodes of $e$ belong to $\ge 2$ different phases}{
        \texttt{is\_gb[$e$]} $\leftarrow$ \texttt{True}\;
        Assign weakened properties: $f_{t0}^e \leftarrow 0.4\,f_{t0}$,
                                     $A_D^e \leftarrow 3\,A_D$\;
    }
}
Assign material properties per node from \texttt{phase} table\;
\end{algorithm}

\vspace{4pt}
The corresponding Python implementation sketch is:

\begin{lstlisting}[style=pythonstyle, caption={Voronoi microstructure generation (core routine)}, label={lst:voronoi}]
import numpy as np
from scipy.spatial import Voronoi, cKDTree

def generate_voronoi_microstructure(domain, n_grains, mineral_props,
                                     mesh_nodes, n_lloyd=30):
    """
    domain       : (Lx, Ly, Lz)  -- domain dimensions [m]
    n_grains     : int            -- target number of grains
    mineral_props: dict with keys = mineral names,
                   values = {'fraction', 'kappa', 'beta', 'E',
                              'mu', 'rho', 'Cp', 'r_avg'}
    mesh_nodes   : (N, 3) array of FEM node coordinates
    Returns: phase_ids (N,), is_gb (M,) for M elements
    """
    minerals = list(mineral_props.keys())
    fractions = np.array([mineral_props[m]['fraction'] for m in minerals])
    fractions /= fractions.sum()        # normalise

    # --- seed points (Poisson disk approximation via random + thinning) ---
    seeds = np.random.rand(n_grains, 3) * domain

    # --- Lloyd relaxation ---
    for _ in range(n_lloyd):
        vor = Voronoi(seeds)
        new_seeds = np.zeros_like(seeds)
        counts    = np.zeros(n_grains)
        for idx, region_id in enumerate(vor.point_region):
            region = vor.regions[region_id]
            if -1 not in region and len(region) > 0:
                verts = vor.vertices[region]
                new_seeds[idx] = verts.mean(axis=0)
                counts[idx] = 1
        mask = counts > 0
        seeds[mask] = new_seeds[mask]   # update converged seeds

    # --- phase assignment ---
    phase_ids_grains = np.random.choice(
        len(minerals), size=n_grains, p=fractions)

    # --- map nodes to nearest grain ---
    tree = cKDTree(seeds)
    _, nearest_grain = tree.query(mesh_nodes)
    node_phases = phase_ids_grains[nearest_grain]

    return node_phases, minerals, mineral_props
\end{lstlisting}

\subsection{Time-step control}

The adaptive time step is bounded by three conditions:
\begin{equation}
    \Delta t \le \min\!\left(
        C_\mathrm{th}\frac{(\Delta x)^2}{\alpha},\;
        C_D\frac{1 - D}{|\dot{D}|},\;
        C_\mathrm{sp}\frac{h_\mathrm{spall}}{\mathrm{RoP}}
    \right),
    \label{eq:cfl}
\end{equation}
where $C_\mathrm{th} \approx 0.4$ is the thermal CFL constant, $C_D \approx 0.1$
limits damage increment per step, and $C_\mathrm{sp}$ prevents a spall event
spanning more than one time step.

\subsection{Adaptive mesh refinement with AMReX/ALAMO}
\label{sec:amr}

The enthalpy PDE \cref{eq:enthalpy_pde} is implemented within the
\textbf{AMReX} block-structured adaptive mesh refinement (AMR) framework
\citep{zhang2019amrex}, using the \textbf{ALAMO} (AMReX-based Low-Mach
Operations) application layer.  AMR is essential for this problem because
the physically important regions — the ablation front, the grain-boundary
damage zone, and the beam footprint — span length scales from micrometres
(grain boundaries) to centimetres (borehole radius), a ratio of $\sim 10^4$.
Resolving this with a uniform mesh would be computationally prohibitive.

\subsubsection{Block-structured AMR}

AMReX uses a \textbf{block-structured} hierarchy of Cartesian grids.
Starting from a coarse base grid that spans the full domain $\Omega$,
regions requiring higher resolution are covered by progressively finer
patches.  The refinement hierarchy has $\ell = 0, 1, \ldots, L$ levels,
where level $\ell+1$ has twice the spatial resolution of level $\ell$ (a
refinement ratio of 2).  The simulation uses \textbf{two layers of
adaptive refinement}, giving a total resolution ratio of $4:1$ between
the coarsest and finest levels.

Refinement is triggered by tagging criteria applied at each time step.
Three tagging conditions are used:
\begin{enumerate}[noitemsep]
    \item \textbf{Enthalpy gradient:} refine cells where
          $|\nabla H| > \tau_H$, placing fine cells at the ablation
          front and phase interfaces.
    \item \textbf{Damage gradient:} refine cells where
          $|\nabla D| > \tau_D$, placing fine cells at the leading edge
          of the damage zone.
    \item \textbf{Beam footprint:} refine cells within the beam spot
          radius $2\omega(z)$, ensuring adequate resolution of the
          Gaussian profile.
\end{enumerate}
Threshold values $\tau_H$ and $\tau_D$ are set as fractions of the maximum
observed gradient at each level, making them automatically scale-adaptive.

\subsubsection{Finite volume discretisation}

The enthalpy equation \cref{eq:enthalpy_pde} is discretised with a
second-order cell-centred finite volume scheme on the AMR hierarchy.
The diffusion flux at cell face $(i+\tfrac{1}{2}, j, k)$ is
\begin{equation}
    F_{i+\frac{1}{2}} = \frac{\kappa_{i+\frac{1}{2}}}
                              {\rho_{i+\frac{1}{2}}\,c_{p,i+\frac{1}{2}}}
                        \frac{H_{i+1} - H_i}{\Delta x},
    \label{eq:fv_flux}
\end{equation}
where face-centred properties are computed by linear interpolation.  At
coarse--fine grid interfaces, flux matching is enforced via \emph{reflux}
operations, which correct the coarse-level fluxes to be consistent with
the fine-level solution, ensuring conservation across levels.

\subsubsection{Time integration}

Each AMR level advances with its own time step $\Delta t^\ell$, subject to
the CFL condition at that level:
\begin{equation}
    \Delta t^\ell \le C_\mathrm{th}
                    \frac{(\Delta x^\ell)^2}
                    {\max(\kappa/(\rho c_p))},
    \label{eq:cfl_amr}
\end{equation}
with $C_\mathrm{th} \approx 0.4$.  The enthalpy equation is advanced
with an \textbf{implicit Crank--Nicolson} scheme at each level, providing
unconditional stability for the diffusion term.  The source terms
$\dot{Q}$ and surface losses are treated explicitly, which is stable
provided $\Delta t^\ell$ satisfies \cref{eq:cfl_amr}.  Subcycling in
time is used: the fine level takes multiple steps per coarse step,
with the ratio equal to the refinement ratio (2 per level).

\subsubsection{Multi-level synchronisation}

After each coarse time step, a \textbf{synchronisation sweep} is
performed:
\begin{enumerate}[noitemsep]
    \item Reflux: correct coarse-level $H$ using fine-level flux registers.
    \item Average-down: overwrite coarse cells covered by fine cells with
          the arithmetic average of fine cell values.
    \item Regrid: re-evaluate tagging criteria and regenerate the AMR
          hierarchy if the number of tagged cells outside existing fine
          patches exceeds a threshold.
\end{enumerate}
Regridding is performed every $N_\mathrm{regrid} = 4$ coarse steps to
balance accuracy against the overhead of grid generation.

\subsubsection{Parallelism}

AMReX decomposes the grid hierarchy into boxes and distributes them
across MPI ranks using a \textbf{space-filling curve} load-balancing
algorithm.  Within each box, OpenMP threading is used for the innermost
loops.  This hybrid MPI/OpenMP approach scales efficiently to $O(10^4)$
cores, enabling full three-dimensional simulations of the borehole at
grain-scale resolution within practical wall-clock times.

\subsubsection{Integration with Voronoi microstructure}

The Voronoi grain structure is stored as a cell-centred integer array
(phase identifier) on the base AMR level.  When the hierarchy is
refined, the phase array is propagated to fine levels by
\emph{injection} (nearest-neighbour copy, preserving sharp grain
boundaries without artificial smearing).  Material properties
$\{\kappa_k, \rho_k, c_{p,k}, E_k, \beta_k\}$ are looked up per cell
from the phase identifier at runtime.

\begin{notebox}
\textbf{Implementation note.}  The enthalpy inversion
(\S\ref{sec:heat}) is performed cell-by-cell immediately after each
implicit solve, using the phase-indicator enthalpies
$\{H^S, H^L, H^{LV}, H^V\}$ stored as level-0 metadata.  Newton
iteration is used only in the single-phase regions where the
$T \to H$ relation is nonlinear due to temperature-dependent $c_p(T)$;
the phase-transition intervals use the closed-form expressions
\cref{eq:phase_fracs} directly.
\end{notebox}

% =============================================================================
\section{Validation}
\label{sec:validation}
% =============================================================================

Two validation cases are selected to test the most critical components of the
model: the thermal field and the spallation onset.

\subsection{Case 1 — MMW temperature field and rate of penetration
            \textnormal{\citep{zhang2023thermal}}}
\label{sec:val_zhang}

\textbf{Setup.}  \citet{zhang2023thermal} present the only available experimental
validation dataset for MMW drilling of geological materials.  They measure
temperature profiles as a function of depth and radial position in Barre granite,
and record the rate of penetration as a function of beam power under vaporisation
conditions.  The experimental parameters are: $P_0 \approx 10$--$50\,\mathrm{kW}$;
beam waist $w_0 \approx 10$--$20\,\mathrm{mm}$; wavelength $\lambda = 3\,\mathrm{mm}$;
Barre granite with $\kappa \approx 2.9\,\mathrm{W\,m^{-1}K^{-1}}$,
$\rho = 2630\,\mathrm{kg\,m^{-3}}$.

\textbf{What is tested.}  This case validates the enthalpy conservation model
(\cref{eq:enthalpy_pde}--\cref{eq:source}), the Gaussian beam implementation
(\cref{eq:intensity}--\cref{eq:waist}), and the vaporisation RoP formula
(\cref{eq:rop_vap}).  The spallation damage model is inactive at these high power
levels.

\textbf{Pass criteria.}
\begin{itemize}[noitemsep]
    \item Simulated temperatures within 15\% of measured values at depths of
          5, 10, 20\,mm below the surface.
    \item Simulated RoP within 20\% of measured values across the power range.
    \item Temperature radial profile (FWHM) matches experimental spot size to
          within 10\%.
\end{itemize}

\subsection{Case 2 — Lowest required surface temperature and onset-spallation
            time \textnormal{\citep{hu2018lrst}}}
\label{sec:val_hu}

\textbf{Setup.}  \citet{hu2018lrst} conducted flame-jet spallation experiments on
$100\times100\times100\,\mathrm{mm}$ granite and sandstone specimens and measured
(i) the LRST by infrared thermometer and (ii) the onset-spallation time by
direct observation.  Heat flux was $Q \approx 6$--$8.5\times10^4\,\mathrm{W\,m}^{-2}$
over a 30\,mm diameter heating zone.  Material properties are given in
\cref{tab:matprops}.

\textbf{What is tested.}  This case validates the thermal--damage coupling:
the heat equation drives the stress field, which drives the damage evolution, which
predicts the onset time.  The Voronoi microstructure is critical here — homogeneous
models are known to underpredict the breakage factor by $\sim 10$--$18\%$.

\textbf{Pass criteria.}
\begin{itemize}[noitemsep]
    \item LRST: $791\,\mathrm{K} \pm 10\%$ (granite); $842\,\mathrm{K} \pm 10\%$
          (sandstone).
    \item Onset time: $31.5\,\mathrm{s} \pm 20\%$ (granite);
          $65.75\,\mathrm{s} \pm 20\%$ (sandstone).
    \item Damage-zone depth: $1$--$5\,\mathrm{mm}$ for granite.
    \item Heterogeneous model should yield $f_b \approx 0.91$ at onset (granite),
          vs.\ $f_b \approx 0.82$ for homogeneous, consistent with experiment.
\end{itemize}

\begin{table}[H]
\centering
\caption{Material properties for validation cases.  Granite and sandstone from
         \citet{hu2018lrst}; mineral phases from \citet{vogler2020simulation}.}
\label{tab:matprops}
\renewcommand{\arraystretch}{1.3}
\begin{tabular}{lllll}
\toprule
\textbf{Property} & \textbf{Symbol} & \textbf{Unit}
    & \textbf{Granite} & \textbf{Sandstone} \\
\midrule
Density                  & $\rho$    & kg\,m$^{-3}$ & 2640   & 2480  \\
Young's modulus          & $E_0$     & GPa          & 29.98  & 16.63 \\
Poisson's ratio          & $\mu$     & ---          & 0.19   & 0.34  \\
Internal friction angle  & $\phi$    & $^\circ$     & 30     & 25    \\
Cohesion                 & $c$       & MPa          & 20     & 12    \\
Thermal conductivity     & $\kappa_0$& W\,m$^{-1}$K$^{-1}$ & 2.50 & 4.96 \\
Specific heat            & $C_p$     & J\,kg$^{-1}$K$^{-1}$ & 827  & 829  \\
Thermal expansion        & $\beta$   & K$^{-1}$     & $8\times10^{-6}$ & $1\times10^{-5}$ \\
LRST (experimental)      & ---       & K            & 791    & 842   \\
Onset time (experimental)& $t_s$     & s            & 31.5   & 65.75 \\
\bottomrule
\end{tabular}
\end{table}

% =============================================================================
\section{Proposed Results}
\label{sec:results}
% =============================================================================

\subsection{Spallation vs vaporisation regime map}

The first result to be generated is a \textbf{regime map} in the
$(P_0,\, z_d)$ plane (beam power vs.\ drilling depth) showing where each removal
mode dominates.  At shallow depth and low power, spallation dominates; at depth
and high power, vaporisation dominates.  The transition boundary is defined by
\begin{equation}
    \mathrm{RoP}_\mathrm{spall}(P_0, z_d, \vv{\theta}) =
    \mathrm{RoP}_\mathrm{vap}(P_0, z_d, \vv{\theta}),
    \label{eq:transition}
\end{equation}
where $\vv{\theta}$ collects rock properties.  The EOS enters through
$E(P_c, T)$: at depth, the raised modulus increases the energy required for
spallation while the vaporisation energy is depth-independent at leading order.
The regime map will be computed for each of the three target rock types.

\textbf{Expected finding.}  We anticipate a monotone transition from spallation
to vaporisation as depth increases, with the transition depth being
material-dependent and strongly correlated with Young's modulus and quartz
content.

\subsection{Material comparison: granite, basalt, and limestone}

Three rock types are selected to span a wide range of mineralogy and mechanical
properties:

\begin{itemize}[noitemsep]
    \item \textbf{Granite}: silica-rich ($\sim 30\%$ quartz), high $E$, high
    $\kappa_\mathrm{qtz}$ mismatch — expected to be the most
    spallation-susceptible rock.
    \item \textbf{Basalt}: fine-grained, low quartz content, high density —
    expected lower spallation susceptibility; vaporisation threshold lower due
    to lower $T_\mathrm{melt}$ ($\approx 1100\,^\circ\mathrm{C}$).
    \item \textbf{Limestone}: calcite-dominated, undergoes $\mathrm{CaCO_3}
    \to \mathrm{CaO} + \mathrm{CO_2}$ decomposition at $\approx 850\,^\circ
    \mathrm{C}$, releasing gas that may contribute an additional pressure-driven
    spallation mechanism not present in silicates.
\end{itemize}

\noindent
For each material the simulation will compute: LRST; onset-spallation time;
RoP as a function of $P_0$; and specific energy (energy per unit volume removed).
Material properties will be taken from the compiled databases of
\citet{vogler2020simulation} and supplemented from \citet{liu2023failure} for
mineral phases not covered therein.

\subsection{Spall initiation time scaling validation}
\label{sec:scaling_validation}

The analytical scaling law \cref{eq:spall_thickness} predicts
\begin{equation}
    t_\mathrm{spall} \propto \frac{K_r^2}{\rho^2\,C_p^2\,Q^2}\,
    \frac{1}{\pi\alpha},
    \label{eq:t_scaling}
\end{equation}
i.e.\ onset time scales as $Q^{-2}$ and $\alpha^{-1}$.  This will be tested
against:
\begin{enumerate}[noitemsep]
    \item Experimental data of \citet{hu2018lrst} for granite and sandstone
          across three heat-flux levels.
    \item Full numerical simulation results from the present model across a
          grid of $(Q, \alpha)$ values.
    \item Literature data of \citet{wilkinson1993experimental} for flame-jet
          spallation of granite at varied stand-off distances (and hence varied
          effective $Q$).
\end{enumerate}
A log-log regression of simulated $t_\mathrm{spall}$ against $Q$ will test the
predicted $-2$ slope; deviations indicate the importance of mechanisms beyond
the simple diffusion scaling (e.g.\ grain-boundary kinetics or quartz transition
effects).

\subsection{Depth dependence and EOS effects}

Using the Mie--Gr\"{u}neisen EOS, the variation of key outputs with drilling
depth will be computed from the surface to $z_d = 10\,\mathrm{km}$:
\begin{itemize}[noitemsep]
    \item $\mathrm{RoP}(z_d)$ for fixed $P_0$, showing the expected decrease
          with depth.
    \item The depth at which the dominant removal mode transitions from
          spallation to vaporisation.
    \item Required beam power to maintain a target $\mathrm{RoP} = 5\,\mathrm{m/h}$
          as a function of depth — a direct engineering design output.
\end{itemize}

\textbf{Expected finding.}  Spallation RoP is expected to fall faster with depth
than vaporisation RoP, because the confining pressure correction to the damage
threshold \cref{eq:depth_correction} scales as $P_c \sim z_d$, while the
vaporisation energy $L_\mathrm{tot}$ is depth-independent.  The transition depth
is expected to be of order $1$--$3\,\mathrm{km}$ for granite under typical
operational beam powers.

% =============================================================================
\section{Conclusions}
\label{sec:conclusions}
% =============================================================================

This proposal describes a unified thermomechanical model for MMW-driven rock
drilling that for the first time treats spallation and vaporisation as two
competing modes within a single engineering framework.  The key modelling
contributions are:

\begin{enumerate}[noitemsep]
    \item A \textbf{Voronoi heterogeneous microstructure} with grain-boundary
    cohesive zones as the physical mechanism driving spallation — without
    which no purely continuum model can predict the experimentally observed
    damage patterns.

    \item A \textbf{Drucker--Prager failure criterion} replacing the
    pressure-independent von Mises criterion, enabling correct prediction of
    the depth-dependent damage threshold.

    \item A \textbf{Mie--Gr\"{u}neisen equation of state} linking elastic
    modulus to pressure and internal energy, capturing the mechanical stiffening
    of rock at depth and its consequences for the spallation--vaporisation
    transition.

    \item A \textbf{continuous damage evolution law} driven by the DP overstress,
    with mineral-phase-specific rate coefficients and grain-boundary weakening,
    replacing the binary threshold criteria of prior work.

    \item A \textbf{spall thickness scaling law} derived from thermal diffusion
    theory, providing a fast engineering estimate of RoP without requiring full
    simulation.
\end{enumerate}

The proposed validation against \citet{zhang2023thermal} and \citet{hu2018lrst}
provides a rigorous test of the thermal and damage sub-models respectively.
The four result directions — regime maps, material comparison, scaling law
validation, and depth dependence — collectively address the core engineering
question: \emph{for a given rock type and target depth, what beam power is
needed to drill at a commercially viable rate?}

% =============================================================================
% Bibliography
% =============================================================================
\bibliographystyle{plainnat}

\begin{thebibliography}{99}

\bibitem[Alexiades \& Solomon(1993)]{alexiades2010enthalpy}
Alexiades, V.\ \& Solomon, A.D.\ (1993).
\textit{Mathematical Modeling of Melting and Freezing Processes}.
Taylor \& Francis, Washington DC.

\bibitem[Augustine(2009)]{augustine2009hydrothermal}
Augustine, C.R.\ (2009).
Hydrothermal spallation drilling and advanced energy conversion technologies
for engineered geothermal systems.
Ph.D.\ thesis, Massachusetts Institute of Technology.

\bibitem[Carpenter et al.(1998)]{carpenter1998quartz}
Carpenter, M.A., Salje, E.K.H., Graeme-Barber, A., Wruck, B., Dove, M.T.,
\& Knight, K.S.\ (1998).
Calibration of excess thermodynamic properties and elastic constant variations
associated with the $\alpha \leftrightarrow \beta$ phase transition in quartz.
\textit{American Mineralogist}, 83, 2--22.

\bibitem[Griffin et al.(1995)]{griffin1995acoustic}
Griffin, P.G., Agrawal, D., \& Pegg, I.L.\ (1995).
$\alpha/\beta$ phase transition in quartz monitored using acoustic emissions.
\textit{Geophysical Journal International}, 120, 775--782.

\bibitem[Haussühl(1960)]{haussuehl1960quartz}
Hauss\"{u}hl, S.\ (1960).
The $\alpha$--$\beta$ transformation of quartz.
\textit{Nature}, 145, 147--148.

\bibitem[Johnson et al.(2021)]{johnson2021quartz}
Johnson, T.E., Brown, M., Gardiner, N.J., Kirkland, C.L., \&
Smithies, R.H.\ (2021).
The quartz $\alpha \leftrightarrow \beta$ phase transition: Does it drive
damage and reaction in continental crust?
\textit{Earth and Planetary Science Letters}, 553, 116628.

\bibitem[Shen et al.(1993)]{shen1993quartz}
Shen, A., Bassett, W.A., \& Chou, I.M.\ (1993).
The $\alpha$--$\beta$ quartz transition at high temperatures and pressures
in a diamond-anvil cell by laser interferometry.
\textit{American Mineralogist}, 78, 694--698.


Augustine, C.R. (2009).
Hydrothermal spallation drilling and advanced energy conversion technologies
for engineered geothermal systems.
Ph.D.\ thesis, Massachusetts Institute of Technology.

\bibitem[Carslaw \& Jaeger(1959)]{carslaw1959}
Carslaw, H.S.\ \& Jaeger, J.C.\ (1959).
\textit{Conduction of Heat in Solids}.
Oxford University Press.

\bibitem[Gaston et al.(2009)]{gaston2009moose}
Gaston, D., Newman, C., Hansen, G., \& Lebrun-Grandi\'{e}, D.\ (2009).
MOOSE: a parallel computational framework for coupled systems of nonlinear
equations.
\textit{Nuclear Engineering and Design}, 239, 1768--1778.

\bibitem[Houde et al.(2021)]{houde2021mmw}
Houde, A.J., et al.\ (2021).
Unlocking deep superhot rock resources through millimetre-wave drilling.
\textit{GRC Transactions}, 45.

\bibitem[Hu et al.(2018)]{hu2018lrst}
Hu, X., Song, X., Liu, Y., Li, G., Shen, Z., Lyu, Z., \& Liu, Q.\ (2018).
Lowest required surface temperature for thermal spallation in granite and
sandstone specimens: Experiments and simulations.
\textit{Rock Mechanics and Rock Engineering}, 52, 1771--1788.

\bibitem[Woskov(2009)]{woskov2009mmw}
Woskov, P.P.\ \& Cohn, D.R.\ (2009).
Annual report 2009: Millimeter wave deep drilling for geothermal energy,
natural gas and oil.
MIT Energy Initiative Technical Report, MITEI-RP-2009-01.

\bibitem[Zhang et al.(2019)]{zhang2019amrex}
Zhang, W., et al.\ (2019).
AMReX: a framework for block-structured adaptive mesh refinement.
\textit{Journal of Open Source Software}, 4(37), 1370.

\bibitem[Kant et al.(2017)]{kant2017theory}
Kant, M.A., Rossi, E., Madonna, C., H\"{o}ser, D., \& von Rohr, P.R.\ (2017).
A theory on thermal spalling of rocks with a focus on thermal spallation drilling.
\textit{Journal of Geophysical Research: Solid Earth}, 122, 1805--1815.

\bibitem[Liu et al.(2023)]{liu2023failure}
Liu, W., Chang, X., \& Zhu, X.\ (2023).
The failure mechanism of thermal spallation in granitic rock.
Preprint, SSRN 4387953.

\bibitem[Oglesby et al.(2014)]{oglesby2014mmw}
Oglesby, K., et al.\ (2014).
White paper: Drilling and fracturing with millimeter-wave directed energy.
Technical Report, GA-EMS.

\bibitem[Rauenzahn \& Tester(1989)]{rauenzahn1989rock}
Rauenzahn, R.M.\ \& Tester, J.W.\ (1989).
Rock failure mechanisms of flame-jet thermal spallation drilling — theory
and experimental testing.
\textit{International Journal of Rock Mechanics and Mining Sciences}, 26, 381--399.

\bibitem[Rossi et al.(2018)]{rossi2018effects}
Rossi, E., Kant, M.A., Madonna, C., Saar, M.O., \& von Rohr, P.R.\ (2018).
The effects of high heating rate and high temperature on the rock strength:
feasibility study of a thermally assisted drilling method.
\textit{Rock Mechanics and Rock Engineering}, 51, 2957--2964.

\bibitem[Tester et al.(2006)]{tester2006future}
Tester, J.W., et al.\ (2006).
\textit{The Future of Geothermal Energy}.
Massachusetts Institute of Technology.

\bibitem[Vogler et al.(2020)]{vogler2020simulation}
Vogler, D., Walsh, S.D.C., von Rohr, P.R., \& Saar, M.O.\ (2020).
Simulation of rock failure modes in thermal spallation drilling.
\textit{Acta Geotechnica}, 15, 2327--2340.

\bibitem[Wilkinson \& Tester(1993)]{wilkinson1993experimental}
Wilkinson, M.A.\ \& Tester, J.W.\ (1993).
Experimental measurement of surface temperatures during flame-jet induced
thermal spallation.
\textit{Rock Mechanics and Rock Engineering}, 26, 29--62.

\bibitem[Zhang et al.(2023)]{zhang2023thermal}
Zhang, A.Z., Millmore, S.T., \& Nikiforakis, N.\ (2023).
Thermal simulation of millimetre wave ablation of geological materials.
\textit{Computers and Geotechnics}, 163, 105732.

\end{thebibliography}

\end{document}